\documentclass[11pt]{article}
\usepackage[margin=24mm]{geometry}
\usepackage[T1]{fontenc}
\usepackage{lmodern,amsmath,amssymb,amsthm,mathtools,booktabs,array}
\usepackage{xcolor,listings,microtype,hyperref}
\hypersetup{colorlinks=true,linkcolor=blue!50!black,urlcolor=blue!50!black,citecolor=blue!50!black,
pdftitle={CMF conservation and an irrationality measure for zeta(2)}}
\lstset{basicstyle=\ttfamily\footnotesize,breaklines=true,columns=fullflexible,
 keepspaces=true,showstringspaces=false,frame=single,rulecolor=\color{black!20}}
\newtheorem{theorem}{Theorem}
\newtheorem{lemma}{Lemma}
\newcommand{\Z}{\mathbb Z}
\newcommand{\Q}{\mathbb Q}
\newcommand{\Fp}{\mathbb F_p}
\newcommand{\lcm}{\operatorname{lcm}}
\newcommand{\med}{\operatorname{med}}
\allowdisplaybreaks
\setlength{\emergencystretch}{2em}
\title{CMF conservation as an arithmetic tool\\
\large A proof of $\mu(\zeta(2))<4.62675$ for the direction\\
$(10,9,11,-11;2,1,3)$}
\author{Mathematical derivation with reproducible exact certificates}
\date{10 October 2026}
\begin{document}
\maketitle
\begin{abstract}
Conservation allows a matrix product to be decomposed into paths and
blocks whose arithmetic is easier to see. Different initial rows and
different primes may favor different certificates for the same endpoint.
We explain this with an explicit $2\times2$ field, then use four long
blocks of the native $\,{}_4F_3(1)$ field to derive an exact finite
readout for the specified trajectory. A finite-field polynomial
factorization proves the additional divisibility in all six required
prime intervals. A Barnes integral identifies the error exactly and
removes the dominant recurrence mode; an elementary cone argument
then gives the sharp exponential upper rate. A nonzero observability
determinant supplies independence at three consecutive depths.
These ingredients give
$\mu(\zeta(2))\le62636575/13537936<4.62675$.
The finite chamber and numerical constant certificates use exact
rational arithmetic. No formal proof assistant is used.
\end{abstract}

\section{The idea: one endpoint, several arithmetic descriptions}
An arithmetic estimate need not use the same path for every initial row
or every prime. Conservation makes the endpoint product independent of
the admissible path, so each valid estimate constrains the same rational
vector. A favorable first block may preserve an eigenline or a congruence
lattice for one row, while another row benefits from another block order.

For the present trajectory, the effective implementation is a fixed
four-block route followed by prime-dependent evaluation of its exact
readout. The last $11N$ moves are treated as one finite-difference
operator. Their cancellation survives reduction modulo a prime and
produces the extra common divisor. Estimating these moves one at a time
would discard the polynomial degree reduction which proves that divisor.


% LaTeX fragment: exact elementary lemmas and a 2-by-2 example.
% Depends only on amsmath/amssymb and an ordinary theorem environment.
\subsection{Conservation changes the available arithmetic certificates}

Use row vectors and multiply on the right.  A conservative matrix field obeys
\[
 M_x(w)M_y(w+e_x)=M_y(w)M_x(w+e_y).
\]
Thus two admissible paths related by elementary square exchanges have the
same endpoint product.  The evaluation points move when the order changes:
this is not permission to commute two matrices evaluated at the same point.
All exchanges must stay within the domain where the required matrices and
identities are defined.

For a prime $\ell$, let $c_\ell(V)$ denote the minimum $\ell$-adic
valuation of the entries of a rational matrix or row $V$, with
$v_\ell(0)=+\infty$.  The elementary valuation inequality gives
\[
 c_\ell(UV)\geq c_\ell(U)+c_\ell(V).
\]
Indeed, every summand in every entry of $UV$ has valuation at least the
right-hand side, and addition cannot decrease that lower bound.

Suppose a path $\gamma$ is decomposed into full blocks
$B_{\gamma,1},\ldots,B_{\gamma,k}$, and let its endpoint product be $W$.
For an initial row $r$, a useful certificate is
\[
 L_{\ell,r}(\gamma)
 =c_\ell(rB_{\gamma,1})+
   \sum_{j=2}^k c_\ell(B_{\gamma,j}),
 \qquad
 c_\ell(rW)\geq L_{\ell,r}(\gamma).
\]
More accurate certificates retain the transported row through later
blocks, or retain a lattice of possible rows.  The blocks can contain any
number of elementary steps.  Replacing a block by separate entrywise
estimates can discard cancellations, since the displayed product
inequality can be strict.

If $\mathcal P$ is any family of admissible paths to the same endpoint,
conservation therefore gives
\[
 c_\ell(rW)\geq\max_{\gamma\in\mathcal P}
 L_{\ell,r}(\gamma).
\]
The maximizing path can depend on $r$, on $\ell$, and on the endpoint
parameter $N$.  These choices are compatible because every certificate
concerns the same rational endpoint vector.  They do not change the
approximant or its actual reduced denominator.

For several initial rows, the useful common-content bound is
\[
 \min_r\max_{\gamma\in\mathcal P}L_{\ell,r}(\gamma)
 \ \geq\
 \max_{\gamma\in\mathcal P}\min_r L_{\ell,r}(\gamma).
\]
The left side allows a different path for each initial row.  The inequality
can be strict even for a $2\times2$ field and one elementary square.

\subsection{An exact toy example with both kinds of dependence}

Consider the rational potential
\[
 F(x,y)=
 \begin{pmatrix}
 1+x+2y+2xy&3y(1-x)\\
 3x(1-y)&1+2x+y+2xy
 \end{pmatrix},
 \qquad
 M_i(w)=F(w)^{-1}F(w+e_i).
\]
This defines a rational conservative field on its nonsingular domain:
both products around a square telescope to
$F(w)^{-1}F(w+e_x+e_y)$.  The four corners used below are nonsingular.
Writing $X$ for an $x$-step and $Y$ for a $y$-step, their matrices are
\[
 A=M_x(0,0)=\begin{pmatrix}2&0\\3&3\end{pmatrix},\quad
 B=M_y(0,0)=\begin{pmatrix}3&3\\0&2\end{pmatrix},
\]
\[
 C=M_y(1,0)=\begin{pmatrix}3&0\\-3&2\end{pmatrix},\quad
 D=M_x(0,1)=\begin{pmatrix}2&-3\\0&3\end{pmatrix}.
\]
Direct multiplication gives
\[
 W=AC=BD=6I_2,
 \qquad AB\ne BA.
\]
The first coordinate row $e_1=(1,0)$ is a left eigenvector of $A$ with
eigenvalue $2$; the second coordinate row $e_2=(0,1)$ is a left eigenvector
of $B$ with the same eigenvalue.  In contrast,
\[
 e_1B=(3,3),\qquad e_2A=(3,3).
\]
All four whole matrices have content zero at both primes $2$ and $3$.
Consequently a bound using only their individual matrix contents certifies
no positive content.  Retaining the first transported row gives instead:
\[
\begin{array}{c|c|cc|c}
\text{initial row}&\ell&L_{\ell,r}(XY)&L_{\ell,r}(YX)&
\text{best order}\\\hline
e_1&2&1&0&XY\\
e_2&2&0&1&YX\\
e_1&3&0&1&YX\\
e_2&3&1&0&XY
\end{array}
\]
For $\ell=2$, send $e_1$ along $XY$ and $e_2$ along $YX$.  This certifies
that both endpoint rows are divisible by $2$.  For $\ell=3$, reverse those
choices; both are divisible by $3$.  Since the endpoint vectors agree
exactly along the two paths, both endpoint rows are divisible by $6$.
Either single order used for both rows gives only the common-content lower
bound zero at each prime under this particular certificate.  The endpoint
calculation $W=6I_2$ independently confirms that the optimized bounds are
sharp.

This example is deliberately small enough that the endpoint product can be
inspected directly.  Its purpose is to exhibit exactly the information
lost by an arithmetic estimate, and how conservation recovers that
information.  For large hypergeometric products one seeks symbolic block
identities that supply the same certificates without expanding the full
endpoint matrix.

\paragraph{The same effect as denominator clearing.}
There is also a rational version of the square with edge matrices
$A/6,B/6,C,D$, so both endpoint products are $I_2$.  It is realized by
bilinearly interpolating potential values
$I_2,A/6,B/6,I_2$ at the four corners and again using
$M_i(w)=F(w)^{-1}F(w+e_i)$.
For $e_1$, the $XY$ calculation clears denominators by $3$, whereas the
$YX$ calculation clears them by $2$.  For $e_2$ the roles reverse.
Since either calculation concerns the same endpoint row, localizing at
each prime certifies integrality: use the path with no $2$-denominator at
$2$ and the path with no $3$-denominator at $3$.  Neither whole-matrix
denominator estimate alone displays this fact.

Equivalently, if the positive integers $D_\gamma$ all clear denominators
of the same rational row, then $\gcd_\gamma D_\gamma$ clears them as well,
by B\'ezout's identity.  For two rows, independently optimizing their
certificates gives the common denominator
\[
 \operatorname{lcm}\!\left(
 \gcd_\gamma D_{1,\gamma},\,
 \gcd_\gamma D_{2,\gamma}\right).
\]
This divides the less flexible expression
\[
 \gcd_\gamma\operatorname{lcm}
 \bigl(D_{1,\gamma},D_{2,\gamma}\bigr).
\]
The rational toy square gives $1$ for the first expression and $6$ for
the second.  Prime-specific ordering is precisely the local form of
taking these greatest common divisors; it does not splice different
approximants together.

\paragraph{What an arithmetic proof must record.}
A usable certificate specifies the initial row, the prime range, the
entire block order, every block's starting point, and the retained
transported row or lattice.  Numerical endpoint readouts suggest such
certificates; a uniform identity or valuation argument must establish
them for all admissible parameters.  A finite set of successful readouts
alone is not that uniform argument.


\section{Comparison with the literature}

% Include in the main article before its bibliography.

Rhin and Viola's permutation-group method compares representations of the
same linear form obtained from transformations of a multiple integral.
Factorial quotients in the transformation formulas supply additional
$p$-adic information; their 1996 paper obtains
$\mu(\zeta(2))<5.441243$ \cite{RhinViola1996}.  Its arithmetic analysis is
particularly explicit in Sections~3--4.  This provides the relevant analogy:
several equivalent descriptions can certify more divisibility than one
description considered in isolation.

The choice of description separately at each prime has a direct precedent.
In Zudilin's hypergeometric treatment, Lemma~7 maximizes valuations of
factorial quotients over parameter permutations.  Section~6 then compares
two different hypergeometric descriptions of the same pair of coefficients
and takes the pointwise maximum of their arithmetic profiles,
$\widetilde\varphi(x)=\max\{\varphi(x),\widehat\varphi(x)\}$.
Equations~(28)--(29), followed by the asymptotic calculation in Section~7,
yield $\mu(\zeta(2))\leq5.09541178\ldots$ \cite{Zudilin2014}.
Thus prime-dependent selection should be presented as a familiar arithmetic
principle, here organized through CMF path identities.

CMFs also already have an arithmetic role in the literature.
David's potential matrix is independent of the lattice path
(Definition~11).  Lemma~36 and Corollary~37 use identities relating
horizontal and vertical directions to establish factorial divisibility
for a $\zeta(3)$ field \cite{David2023}.  The July~2026 revision of
Weinbaum et al. studies CMF asymptotics through a discrete
Levinson-type framework and explicitly separates structural results from
experimental arithmetic conjectures \cite{Weinbaum2026}.

In the present proposal, the representations are shifted matrix products
along admissible paths in a known conserving family.  A selected row is
transported exactly into each block, so a certificate can retain congruence
information lost by an entrywise matrix estimate.  One may then optimize the
certificate for each row and prime, while conservation keeps the endpoint
readout fixed.  This is the concrete CMF mechanism to investigate; the extra prime intervals for the present $\zeta(2)$ trajectory are proved by the uniform polynomial argument in Appendix~\ref{app:all-shells}.  Suitable descriptions of the classical work are
``technically intricate'', ``case-intensive'', and
``dependent on carefully chosen transformations''.



\section{The specified $\zeta(2)$ trajectory}
We use the native field $\mathrm{pFq}(4,3,1)$, with initial vertex and direction
\[
 w_0=(1,1,1,0;1,1,1),\qquad v=(10,9,11,-11;2,1,3).
\]
Let $W(N)$ be its unscaled endpoint product, acting on rows on the right,
and set
\[
 R=\begin{pmatrix}0&2&-1\\2&0&0\end{pmatrix},\qquad
 \binom{p_N^{\rm raw}}{q_N^{\rm raw}}=RW(N)e_0.
\]
Thus $p_N^{\rm raw}/q_N^{\rm raw}$ is exactly the expression in the
experiment, $(2W_{10}-W_{20})/(2W_{00})$.
The column labels below are $r=0,1,2$.
The two initial rows of $R$ are transported separately when deriving
arithmetic bounds; they are not reset at the boundaries of blocks.

At a vertex $w=(x_0,x_1,x_2,-m;y_0,y_1,y_2)$, write
\[
 Q_w(s)=s\prod_{i=0}^2(s+y_i-1)-(s-m)\prod_{i=0}^2(s+x_i)
       =\sum_{j=0}^3a_js^j,
\quad
 T_w=\begin{pmatrix}0&0&-a_0/a_3\\1&0&-a_1/a_3\\0&1&-a_2/a_3\end{pmatrix}.
\]
An upper increase has matrix $I+T_w/x_i$. A lower increase is the
inverse of $I+T_{w+e_{y_i}}/y_i$, evaluated at the target in $T$.
A decrease of $x_3$ is likewise the inverse upper move evaluated at
its target. These evaluation points are part of the multiplication order.

\subsection{The four chosen blocks}
For $t=3$ or $4$, the following blocks reach the same endpoint:
\[
\begin{array}{c|c|c}
\text{block}&\text{direction per block step}&\text{number of steps}\\\hline
A&(1,1,1,0;0,0,0)&tN\\
B&(0,0,0,0;2,1,3)&N\\
C&(10-t,9-t,11-t,0;0,0,0)&N\\
D&(0,0,0,-1;0,0,0)&11N
\end{array}
\]
Within each step, make its coordinate moves from left to right, always
evaluating at the current vertex. The first block raises all positive
uppers enough that the lower moves are regular. The lower block then
establishes the harmonic shifts; the remaining upper block establishes
the binomial widths; the final negative-upper block performs a long
finite difference. In particular this decomposition treats the last
$11N$ steps together, rather than estimating them independently.

At the endpoint all three paired widths $x_i-y_i$ equal $8N$.
This is the reason that the readout below contains three binomial
polynomials of a common width. Appendix~\ref{app:blocks} records the
complete $t=3$ vertices and compatible integral normalizations, including
every block boundary. Appendix~\ref{app:finite} verifies all generators
and the regularity of both routes uniformly in $N$.

\subsection{Reading the arithmetic from the rows}
Put $d_n=\lcm(1,\ldots,n)$, with $d_0=1$, and let
$H_n^{(2)}=\sum_{k=1}^n k^{-2}$, with $H_0^{(2)}=0$.
Define
\[
\begin{split}
 C_N&=\binom{10N}{2N}\binom{9N}{N}\binom{11N}{3N},\qquad m=11N,\\
 F(s)&=\binom{s+10N}{8N}\binom{s+9N}{8N}\binom{s+11N}{8N},\qquad
 c_j=(-1)^j\binom mjF(j),\\
 S_r(s)&=\sum_{j=0}^m(-1)^j\binom mj
       \frac{s^rF(s)-j^rF(j)}{s-j},\\
 J_\eta(P)&=\sum_{h=0}^{\deg P}\frac{(-1)^h}{h+1}\Delta^hP(-\eta),
 \qquad \Delta P(s)=P(s+1)-P(s).
\end{split}
\]
We use $0^0=1$ and the polynomial value of each divided difference
at a removable singularity. The exact native readout is
\begin{equation}\label{eq:finite}
 C_NRW(N)e_r=
 2\binom{\displaystyle\sum_{j=0}^m j^rc_jH_{j+\eta}^{(2)}+J_\eta(S_r)}
        {\displaystyle\sum_{j=0}^m j^rc_j},
 \qquad 2N\le\eta\le9N.
\end{equation}
This is a uniform identity, proved in Appendix~\ref{app:finite}, and
its value is independent of the admissible integer $\eta$.
The first row contains the harmonic and polynomial correction terms;
the second is an integer binomial sum after multiplication by $C_N$.
The polynomial degree drops from $24N+r-1$ to at most $13N+r-1$
because a finite difference of order $11N$ kills all lower powers.
That is a whole-block cancellation.

The matrix readout therefore does more than suggest a growth rate.
Its exact expression supplies the integer-valued polynomial basis,
the possible denominator indices, and the binomial carries that can
force common divisors. A raw list of rational approximants would hide
this structure. Discovering a conserving family provides a systematic
route to the identities that expose it.



\section{The arithmetic normalization}
For real $x,y$, set
\[
\begin{split}
 \rho(x,y)={}&\lfloor11x\rfloor-\lfloor y\rfloor-\lfloor11x-y\rfloor\\
 &+\lfloor10x+y\rfloor+\lfloor9x+y\rfloor+\lfloor11x+y\rfloor
 -3\lfloor8x\rfloor\\
 &-\lfloor2x+y\rfloor-\lfloor x+y\rfloor-\lfloor3x+y\rfloor,\\
 \phi(x)={}&\min_{0\le y<1}\rho(x,y),\qquad
 G_N=\prod_{\substack{\sqrt{50N}<p\le13N\\p\ \text{prime}}}p^{\phi(N/p)},\qquad
 D_{N,r}=d_{12N}d_{13N+r}.
\end{split}
\]
Each summand in the carry interpretation of $\rho$ is nonnegative.
The large-prime cutoff ensures that only the first level of Legendre's
factorial valuation formula is involved.

Appendix~\ref{app:baseline} proves the baseline statement
\begin{equation}\label{eq:baseline}
 \frac{C_ND_{N,r}}{G_N}RW(N)e_r\in\Z^2
 \quad(N\ge1,\ r=0,1,2).
\end{equation}
It does so by bounding the divided differences in an integer-valued
binomial basis, accounting for a missing harmonic prime by a binomial
carry, and transferring common divisibility through Newton coefficients.
Thus this is a proof of the arithmetic of the actual native rows,
not an extrapolation from the checked endpoints.

The additional common divisor comes from the following six prime intervals.
Let $H_{N,r}$ be the product of primes $p>\sqrt{50N}$ belonging to
one of the following \emph{open} intervals:
\begin{equation}\label{eq:shells}
\begin{gathered}
 ((13N+r)/2,8N),\quad ((13N+r)/3,9N/2),\quad
 ((13N+r)/4,10N/3),\\
 ((13N+r)/5,8N/3),\quad ((13N+r)/6,9N/4),\quad
 ((13N+r)/12,9N/8).
\end{gathered}
\end{equation}

\begin{theorem}[Arithmetic of the native readouts]\label{thm:arithmetic}
For every $N\ge1$ and $r=0,1,2$,
\[
 \boxed{\frac{C_Nd_{12N}d_{13N+r}}{G_NH_{N,r}}RW(N)e_r\in\Z^2.}
\]
\end{theorem}
\begin{proof}
The baseline argument in Appendix~\ref{app:baseline} clears all
denominators and divides out $G_N$. It remains to establish one further
factor of each prime in the six open intervals. Appendix~\ref{app:all-shells}
proves this simultaneously for all six intervals and all three columns.
Here is the mechanism, before its detailed calculation.

Write $8N=ap+b$, $11N=Mp+d$, with $0\le b,d<p$, and let
$h=\phi(N/p)$. In $\mathbb F_p[X]$ form
\[
 \mathcal Q(X)=
 \frac{\displaystyle\prod_{u=9}^{11}\binom{X+uN}{b}}
 {(X^p-X)^{h-1}\displaystyle\prod_{j=0}^d(X-j)}.
\]
For $h=0$, the expression means multiplication by $X^p-X$ in the
numerator. The carry lower bound proves that this is a polynomial:
every denominator root occurs in the numerator with at least the
required multiplicity. In a shell labeled by $K=2,3,4,5,6,12$,
\[
 \deg(X^r\mathcal Q)=13N+r-(K-1)p-1<p-1.
\]
Therefore $\sum_{k\in\mathbb F_p}k^r\mathcal Q(k)=0$.

The whole-block rational kernel, reduced at $s=k+pV$, separates into
$\mathcal Q(k)$ times a rational function of the high digit $V$.
The floor-cell table shows that the latter function is independent
of $k$ on the support of $\mathcal Q$. Choose the admissible shift
$\eta=p$ or $2p$. Both the leading harmonic residue and the leading
Newton-coefficient residue then contain the same vanishing sum.
Both parts of the numerator gain a factor of $p$, and the denominator
already has the required factor. Combining the distinct primes proves
the assertion. The appendix accounts for derivative nodes and for the
carry chambers where a termwise argument would fail.
\end{proof}

This identifies precisely how the conserving family helps the proof.
It transports the original rows to a useful four-block description;
the exact readout records a long finite difference; and its polynomial
structure supplies the uniform arithmetic cancellation. The choice of
$\eta$ here is a prime-dependent representation of that same readout.
The toy example separately illustrates literal prime-dependent path
orders. We do not identify these two freedoms without an additional
path identity.



\section{The exponential cost of the arithmetic normalization}
Define the arithmetic rate
\[
 \gamma=\int_{1/13}^{\infty}\frac{\phi(x)}{x^2}\,dx.
\]
The function $\phi$ is bounded and one-periodic, and vanishes on
$(0,1/13)$. The prime number theorem, applied on its continuity
intervals, gives $N^{-1}\log G_N\longrightarrow\gamma$.
For completeness, first restrict to $p\ge N/K$, where there are only
finitely many intervals, then bound the omitted part by
$\|\phi\|_\infty\vartheta(N/K)/N=O(1/K)$. The excluded primes
$p\le\sqrt{50N}$ contribute $o(N)$ and do not affect the limit.
Here $\vartheta(x)=\sum_{p\le x}\log p$.

The six interval lengths in \eqref{eq:shells}, divided by $N$, sum to
\[
 \left(8-\frac{13}2\right)+\left(\frac92-\frac{13}3\right)
 +\left(\frac{10}3-\frac{13}4\right)
 +\left(\frac83-\frac{13}5\right)
 +\left(\frac94-\frac{13}6\right)
 +\left(\frac98-\frac{13}{12}\right)=\frac{233}{120}.
\]
Consequently Theorem~\ref{thm:arithmetic} gives the arithmetic
exponent
\[
 c=25-\gamma-\frac{233}{120}<10.929906.
\]
The strict numerical comparison uses a directed rational enclosure,
not a floating-point integral; Appendix~\ref{app:rate} explains it.


In particular, with
$k_N=d_{12N}d_{13N}/(G_NH_{N,0})$, we have
\[
 \lim_{N\to\infty}\frac{\log k_N}{N}
 =25-\gamma-\frac{233}{120}<10.929906.
\]
Only a subexponential factor is lost by omitting primes below
$\sqrt{50N}$ or by using open interval endpoints.


% Analytic completion for the exact native trajectory (10,9,11,-11;2,1,3).
% Depends on the finite formula established in finite_formula_audit.tex.
\section{Analytic estimates and nonvanishing from the native CMF}

Let $W_N$ be the native endpoint product, and write
\[
 \mathsf P_{N,r}=(RW_N)_{0r},\qquad
 \mathsf Q_{N,r}=(RW_N)_{1r},\qquad
 \mathsf E_{N,r}=\zeta(2)\mathsf Q_{N,r}-\mathsf P_{N,r}
 \quad(0\le r\le2).
\]
These symbols denote the rational matrix readouts, before integer
normalization.  Put
\[
 F_N(s)=\prod_{u=9}^{11}\binom{s+uN}{8N},\quad
 C_N=F_N(0),\quad m=11N,\quad
 L_m(s)=\frac{(-1)^m m!}{s(s-1)\cdots(s-m)}.
\]
The finite formula provides the following exact integral identification,
not merely a numerically identified limit.

\subsection{The Barnes integral identifies the error exactly}
For $\eta\ge0$ and a function for which the integral converges, set
\[
 \mathcal I_\eta(g)=\frac1{2\pi i}
 \int_{-\eta-1/2-i\infty}^{-\eta-1/2+i\infty}
 \left(\frac\pi{\sin\pi s}\right)^2g(s)\,ds.
\]
For every polynomial $P$,
\[
 \mathcal I_\eta(P)=J_\eta(P).
\]
Here is a direct verification.  Moving the line one unit to the right
crosses one double pole, whose residue is $P'(-\eta)$; the horizontal
segments tend to zero because of the exponential decay of the kernel.
Thus the integral functional obeys
$\mathcal I_\eta(P(s+1)-P(s))=P'(-\eta)$.
This determines a linear functional on all polynomials, since the
forward-difference map on polynomials is onto.  It agrees with the
binomial-basis identity defining $J_\eta$.

For $j\ge0$, closing a large half-integer rectangle to the right gives
\[
 \mathcal I_\eta\!\left(\frac1{s-j}\right)
 =H_{j+\eta}^{(2)}-\zeta(2).
\]
Indeed, the residue at an integer $k\ne j$ is $-(k-j)^{-2}$, and the
residue at $k=j$ is $\pi^2/3=2\zeta(2)$.  The integral on the distant
right edge tends to zero.  The left edge is oriented upward, which
accounts for the minus sign in the residue sum.

The exact partial-fraction identity
\[
 s^rF_N(s)L_m(s)=S_r(s)+
 \sum_{j=0}^m\frac{j^r(-1)^j\binom mjF_N(j)}{s-j}
\]
therefore yields
\[
 \boxed{\quad C_N\mathsf E_{N,r}
       =-2\mathcal I_\eta(s^rF_NL_m).\quad}                 \tag{A1}
\]
Thus the constant in the error is exactly $\zeta(2)$.
Take $\eta=6N$.  The line may be moved from $\Re s=-6N-1/2$ to
$\Re s=-6N$: the only integer crossed is a zero of $F_N$ of
multiplicity three, so it cancels the double pole of the sine kernel.

\subsection{A coarse exponential bound removes the dominant mode}
Use $t=-s$ and the reflection formula for the gamma function.  Since
\[
 L_m(-t)=-B(t,m+1),
\]
and the three complementary shifts are $N,2N,3N$, whose sum is even,
(A1) becomes
\[
 \frac{C_N\mathsf E_{N,r}}2
 =\frac1{2\pi i}\int_{6N-i\infty}^{6N+i\infty}
 (-t)^r B(t,11N+1)
 \frac{\prod_{v=1}^3\Gamma(t-vN)\Gamma((v+8)N+1-t)}
      {\Gamma(8N+1)^3}
 \frac{\sin\pi t}{\pi}\,dt.                              \tag{A2}
\]
All gamma arguments in this formula have positive real part along the
line.  For $a>0$, real $y$, and $0<\theta<\pi/2$, rotation of the Euler
integral by $\operatorname{sgn}(y)\theta$ proves
\[
 |\Gamma(a+iy)|\le\Gamma(a)(\cos\theta)^{-a}
                                  e^{-\theta|y|}.
\]
The rotation is justified by exponential decay on the outer arc and
$a>0$ on the inner arc.  Also, directly from the beta integral,
$|B(6N+iy,11N+1)|\le B(6N,11N+1)$.

The six gamma arguments in (A2), at $t=6N+iy$, have real parts
$3N,3N+1,4N,4N+1,5N,5N+1$.  Choose $\theta=\pi/4$.
Their combined bound supplies
$2^{12N+3/2}e^{-3\pi|y|/2}$; the sine contributes at most
$e^{\pi|y|}/2$.  Consequently the remaining integral is bounded by
a constant times $N^r$, uniformly for $N\ge1$.  More explicitly,
\[
 |C_N\mathsf E_{N,r}|
 \le K_r N^r 2^{12N}
 B(6N,11N+1)
 \frac{\prod_{k=3}^5\Gamma(kN)\Gamma(kN+1)}
      {\Gamma(8N+1)^3},                                  \tag{A3}
\]
where $K_r$ is independent of $N$.  Stirling's formula gives
\[
 \limsup_{N\to\infty}\frac1N\log|C_N\mathsf E_{N,r}|
 \le h,
\]
\[
 h=12\log2+6\log6+11\log11-17\log17
       +2\sum_{k=3}^5 k\log k-24\log8
   <-18.                                                 \tag{A4}
\]
Since $C_N\ge1$, this already proves exponential decay of the entire
unscaled error row.  It is deliberately a coarse bound; it will identify
which invariant modes are present, after which the matrix recurrence
supplies the sharp upper rate.

\subsection{The limiting matrix and an elementary cone argument}
Let $M_N$ be the native product of the 47 unit moves from the endpoint
at $N$ to that at $N+1$.  Conservation gives $W_{N+1}=W_NM_N$.
For $N\ge1$ put
\[
 S_N=\operatorname{diag}(1,N^{-1},N^{-2}),\qquad
 A_N=S_N^{-1}M_NS_{N+1}.
\]
Every entry is a rational function of $N$, and direct degree comparison
in the native companion matrix gives $A_N=A_\infty+O(N^{-1})$, where
\[
 T=\begin{pmatrix}0&0&10890/13\\1&0&2305/13\\0&1&42/13\end{pmatrix},
\]
\[
 A_\infty=
 \prod_{u\in\{10,9,11\}}(I+T/u)^u
 (I-T/11)^{-11}
 \prod_{v\in\{2,1,3\}}(I+T/v)^{-v}.                       \tag{A5}
\]
To see the limit directly, each unit move changes the vertex by only a
bounded amount, while $S_N^{-1}T_wS_N/N\to T$ and its parameter
denominator divided by $N$ tends to its displayed slope.  All the
limiting factors are invertible.  Since they are rational functions
of the same matrix $T$, their order in (A5) is immaterial.

Write the characteristic polynomial as
$z^3+c_2z^2+c_1z+c_0$.  Exact rational arithmetic gives
\[
 c_0=\frac{1855425871872}{895833302668638834554683750087890625},
\]
\[
\begin{gathered}
 c_1=\frac{6119917581009155345396562793851904}
 {32000061404626447809127858236889541015625},\\
 c_2=\frac{266957778538934989798238376160038092809}
 {15342324857928574809549553125}.
\end{gathered}
\]
Its discriminant is strictly negative.  Thus there is one real
root $-\Lambda$ and a nonreal conjugate pair with common modulus
$\rho$, where
\[
 17400086428<\Lambda<17400086429,\qquad
 \rho^2=\frac{c_0}{\Lambda}<1.                             \tag{A6}
\]
In particular, all three roots are distinct, and $A_\infty$ is
complex-diagonalizable.

For completeness, the needed recurrence conclusion has an elementary
proof.  In a fixed eigenbasis split a row solution of
$x_{N+1}=x_NA_N$ into its dominant coordinate $u_N$ and remaining pair
$v_N$, using the maximum norm on the pair.  For any sufficiently small
$\varepsilon>0$, the perturbation from the diagonal limiting matrix
has norm at most $\varepsilon$ for all sufficiently large $N$.
The cone $|u_N|\ge\|v_N\|$ is then forward invariant, and inside it
\[
 |u_{N+1}|\ge(\Lambda-2\varepsilon)|u_N|,
 \qquad
 \|v_{N+1}\|\le(\rho+2\varepsilon)|u_N|.
\]
Choose $\Lambda-2\varepsilon>\max(1,\rho+2\varepsilon)$.
A nonzero solution tending to zero cannot enter this cone.  Outside
it, $|u_N|<\|v_N\|$ at every sufficiently large index, so
\[
 \|v_{N+1}\|\le(\rho+2\varepsilon)\|v_N\|.
\]
Letting $\varepsilon$ decrease to zero proves a growth upper rate
$\log\rho$ for every decaying solution.  Without the decay condition,
the ordinary operator-norm bound gives upper rate $\log\Lambda$.
The zero solution satisfies both upper bounds automatically.

Apply this first to the fixed error row
$(\zeta(2)R_{1,*}-R_{0,*})W_NS_N$, which tends to zero by (A3), and
then to either fixed row of $RW_NS_N$.  Removing $S_N$ changes only
polynomial factors.  We obtain
\[
 \limsup_N\frac1N\log|\mathsf E_{N,r}|\le\log\rho,
 \qquad
 \limsup_N\frac1N\log|\mathsf Q_{N,r}|\le\log\Lambda.
                                                               \tag{A7}
\]
These are uniform exponential upper estimates, and do not require a
lower estimate for an oscillatory error.

\subsection{Certified numerical rates}
Stirling's formula gives
\[
 \lim_N\frac{\log C_N}{N}
 =c_C=10\log10+9\log9+11\log11
          -2\log2-3\log3-24\log8.
\]
The accompanying standard-library program verifies by exact rational
arithmetic that
\[
 c_C+\log\Lambda<38.168733,
 \qquad c_C+\log\rho<-24.467842.                           \tag{A8}
\]
For reproducibility, logarithms are enclosed by scaling a positive
rational argument into $[1,2]$ and using
\[
 \log x=2\sum_{k=0}^{K-1}\frac{z^{2k+1}}{2k+1}+R_K,
 \quad z=\frac{x-1}{x+1},\quad
 0\le R_K\le\frac{2z^{2K+1}}{(2K+1)(1-z^2)}.
\]
The implementation uses $K=32$, handles signs of all coefficients
when combining intervals, checks the cubic discriminant and the root
bracket (A6), and uses $\log\rho=(\log c_0-\log\Lambda)/2$.
No floating-point value is used in its assertions.

\subsection{A bounded gap of at most two for column zero}
The product $W_N$ is invertible on the chosen regular paths.  The exact
observability determinant is
\[
\begin{gathered}
 \det(e_0,A_\infty e_0,A_\infty^2e_0)
 \\
 =\frac{72143259339220303025804927524593090645506084475442025307751136603648}
 {28469599455181793682074031322586982029718987334183867584429931640625}
 \ne0.
\end{gathered}                                               \tag{A9}
\]
Thus $(e_0,A_Ne_0,A_NA_{N+1}e_0)$ is invertible for all sufficiently
large $N$.  Given integers $a,b$ with $b>0$, the row
\[
 (bR_{0,*}-aR_{1,*})W_NS_N=(-2a,2b,-b)W_NS_N
\]
is nonzero.  It cannot annihilate all three columns in that matrix.
Therefore at least one of
\[
 b\mathsf P_{N,0}-a\mathsf Q_{N,0},\quad
 b\mathsf P_{N+1,0}-a\mathsf Q_{N+1,0},\quad
 b\mathsf P_{N+2,0}-a\mathsf Q_{N+2,0}
\]
is nonzero.  This supplies the needed nonvanishing with a bounded
gap, directly from the native three-dimensional matrices.  Multiplying
each pair by any nonzero normalization preserves this property.

\subsection{The irrationality-measure conclusion}
Theorem~\ref{thm:arithmetic} and the preceding rate calculation supply nonzero scalars $k_N$ with
\[
 P_N=C_Nk_N\mathsf P_{N,0}\in\mathbb Z,\qquad
 Q_N=C_Nk_N\mathsf Q_{N,0}\in\mathbb Z,\qquad
 \limsup_N\frac{\log|k_N|}{N}\le10.929906.
                                                               \tag{A10}
\]
Here $k_N=d_{12N}d_{13N}/(G_NH_{N,0})$.
Equations (A7)--(A8) then give
\[
 \limsup_N\frac{\log|Q_N|}{N}\le\sigma=49.098639,
 \qquad
 \limsup_N\frac{\log|Q_N\zeta(2)-P_N|}{N}
 \le-\tau=-13.537936.
\]
Here is the irrationality-measure step explicitly.  For any small
$\varepsilon>0$, take $n$ just large enough that
$b e^{-(\tau-\varepsilon)n}<1/2$.  At one of $N=n,n+1,n+2$, (A9)
ensures the integer $bP_N-aQ_N$ is nonzero.  Hence
\[
 1\le |bP_N-aQ_N|
 \le b|Q_N|\left|\zeta(2)-\frac ab\right|
          +b|Q_N\zeta(2)-P_N|,
\]
which gives a positive constant times
$b^{-1-(\sigma+\varepsilon)/(\tau-\varepsilon)}$ as a lower bound
for the rational-approximation error.  Letting $\varepsilon\downarrow0$
proves
\[
 \boxed{\quad\mu(\zeta(2))\le1+\frac\sigma\tau
   =\frac{62636575}{13537936}<4.62675.\quad}
\]
All ingredients have now been supplied: the finite CMF identity, the uniform arithmetic normalization, the analytic upper estimates, and bounded-gap independence. Here $\mu(\alpha)$ is the usual irrationality measure: the infimum of $\mu$ such that for every $\varepsilon>0$ only finitely many reduced $a/b$ satisfy $|\alpha-a/b|<b^{-\mu-\varepsilon}$.



\section{The role of the exact code}
The proof reduces its finite checks to explicit rational identities and
finite arrangements. The accompanying Python package verifies the
390 floor cells in the eight-row chamber table, computes the rational
limiting matrix and its discriminant, brackets its real eigenvalue,
checks the observability determinant, and encloses the logarithms with
positive series and explicit tails. Its carry-integral computation also
uses a finite rational partition and a proved tail bound.
These finite certificates check the stated reductions for all parameters
covered by each cell; they are not tests at a chosen finite list of depths.

Other scripts provide additional diagnostics: RamanujanTools checks the
native matrices against the four-block products, the finite-functional
checker checks general vertices, and independent integer arithmetic
checks the local kernel formulas. These diagnostics help catch mistakes
but are not the logical justification for universal statements.

The uniform arithmetic argument keeps the cancellation of the entire
rational kernel. Each finite certificate is used only after its
reduction to a finite set of rational identities or floor cells has
been established in the text.


\appendix
\section{Uniform derivation of the finite readout}\label{app:finite}

% Independently audited finite formula. Insert after the CMF generators are defined.
% All claims below are algebraic; no irrationality-measure conclusion is asserted.


We use row vectors and the unscaled native generators.  Assume
\[
 x_i,y_i\in\mathbb Z,\quad x_i\ge y_i\ge1\quad(0\le i\le2),
 \qquad m\in\mathbb Z_{\ge0},
\]
and put $b_i=x_i-y_i$ and
\[
 I_w=\{\eta\in\mathbb Z:
 \operatorname{med}(y_0-1,y_1-1,y_2-1)
 \le\eta\le\min_i(x_i-1)\}.
\]
Assume $I_w\ne\varnothing$. Define
\[
 F(s)=\prod_{i=0}^2\binom{s+x_i-1}{b_i},\quad C=F(0)>0,
 \quad \omega_j=(-1)^j\binom mj,\quad c_j=\omega_jF(j),
\]
\[
 A(s)=s\prod_i(s+y_i-1),\qquad
 D(s)=(s-m)\prod_i(s+x_i),\qquad Q(s)=A(s)-D(s).
\]
The two monic quartic terms cancel.  Write
$Q(s)=a_0+a_1s+a_2s^2+a_3s^3$ and assume $a_3\ne0$.
For a polynomial $f$, set
\[
 S_f(s)=\sum_{j=0}^m\omega_j
   \frac{f(s)F(s)-f(j)F(j)}{s-j},\qquad
 \ell_\eta(f)=\binom{
 \sum_{j=0}^m f(j)c_jH_{j+\eta}^{(2)}+J_\eta(S_f)}{
 \sum_{j=0}^m f(j)c_j},
\]
where $J_\eta(\binom{s+\eta}{h})=(-1)^h/(h+1)$.
All the divided differences defining $S_f$ are polynomials.

\paragraph{1. The two identities for the finite functional.}
Let $P$ be any polynomial.  Pascal's identity and
\[
 \left.\frac{d}{dt}\binom th\right|_{t=0}
 =\frac{(-1)^{h-1}}h\quad(h\ge1)
\]
give, by checking the binomial basis,
\[
 J_\eta(P(s+1)-P(s))=P'(-\eta).
\]
Translation of that basis also gives
$J_{\eta+1}(P)=J_\eta(P(s-1))$. Consequently
\[
 J_{\eta+1}(P)-J_\eta(P)=-P'(-\eta-1).
\]
These signs fix both cancellation arguments below.

\paragraph{2. The exact rational kernel and the shift of $S_A$.}
As an identity of rational functions,
\[
 L_m(s):=\sum_{j=0}^m\frac{\omega_j}{s-j}
  =\frac{(-1)^m m!}{s(s-1)\cdots(s-m)}.
\]
Indeed, the right side has residue $\omega_j$ at $s=j$, and the
remaining polynomial part vanishes at infinity.  The factor $(-1)^m$
must be retained.  The polynomial $F$ satisfies the rational identity
\[
 \frac{F(s+1)}{F(s)}=
   \frac{\prod_i(s+x_i)}{\prod_i(s+y_i)},\qquad
 \frac{L_m(s+1)}{L_m(s)}=\frac{s-m}{s+1}.
\]
Therefore
\[
 A(s+1)F(s+1)L_m(s+1)=D(s)F(s)L_m(s).
\]
With $u_j=A(j)c_j$, extended by $u_0=u_{m+1}=0$, the same identities
at nonnegative integer arguments give
\[
 D(j)c_j=u_{j+1}\qquad(0\le j\le m).
\]
The endpoint $j=m$ follows from $D(m)=0$; the other cases use the
binomial-weight recurrence.  Since
\[
 S_f(s)=f(s)F(s)L_m(s)-\sum_{j=0}^m\frac{f(j)c_j}{s-j},
\]
the two preceding equations imply, with every residue and endpoint
accounted for,
\[
 \boxed{S_D(s)=S_A(s+1).}
\]
This is a rational-function calculation whose two sides are polynomials;
no division at an individual root of $F$ is being asserted.

\paragraph{3. The companion relation.}
The rational function $A(s)F(s)L_m(s)$ can be written
\[
 \frac{(-1)^m m!}{\prod_i b_i!}
 \frac{\prod_i(s+y_i-1)_{b_i+1}}
 {(s-1)\cdots(s-m)}.
\]
Its denominator is nonzero at $s=-\eta$, including $\eta=0$;
the factor $s$ has already been cancelled.  Since
$\eta\ge\operatorname{med}(y_i-1)$ and
$\eta\le\min_i(x_i-1)$, at least two numerator factors vanish there.
Hence $(AF L_m)'(-\eta)=0$, and differentiation of the expression for
$S_A$ gives
\[
 S_A'(-\eta)=\sum_{j=1}^m\frac{u_j}{(j+\eta)^2}.
\]
On the other hand,
\[
 \sum_{j=0}^m(A(j)-D(j))c_jH_{j+\eta}^{(2)}
 =\sum_{j=1}^m\frac{u_j}{(j+\eta)^2},
\]
by finite telescoping.  The polynomial contribution is
$J_\eta(S_A-S_D)=-S_A'(-\eta)$, and the second component also
telescopes.  Thus
\[
 \boxed{\ell_\eta(Q)=0.}
\]
Only this degree-three relation is needed when multiplying the basis
$1,s,s^2$ by $s$.  It is not necessary, and has not been claimed, that
$\ell_\eta(fQ)=0$ for every polynomial $f$.

\paragraph{4. Independence of the admissible shift.}
If both $\eta$ and $\eta+1$ lie in $I_w$, subtracting the first
components yields
\[
 \ell_{\eta+1}(f)_1-\ell_\eta(f)_1
 =\sum_j\frac{f(j)c_j}{(j+\eta+1)^2}
   -S_f'(-\eta-1)
 =-(fF L_m)'(-\eta-1).
\]
At $s=-\eta-1$, the function $L_m$ has no pole.  At least two of the
factors of $F$ vanish: their integer roots cover the intervals
$y_i,\ldots,x_i-1$, and
$\eta+1\ge\operatorname{med}(y_i)$ and
$\eta+1\le\min_i(x_i-1)$.  The derivative is therefore zero.
The second component contains no $\eta$.  It follows that
$\ell_\eta(f)$ is independent of $\eta\in I_w$.

\paragraph{5. The three types of contiguous move.}
Write
\[
 T_w=\begin{pmatrix}0&0&-a_0/a_3\\1&0&-a_1/a_3\\0&1&-a_2/a_3\end{pmatrix},
 \qquad
 Y_w=\frac2C\bigl(\ell_\eta(1),\ell_\eta(s),\ell_\eta(s^2)\bigr).
\]
The companion relation proves that right multiplication by $T_w$
replaces the three columns by
$2\ell_\eta(s)/C,2\ell_\eta(s^2)/C,2\ell_\eta(s^3)/C$.
For an upper increase $x_i\mapsto x_i+1$,
\[
 F_{\rm new}(s)=\frac{s+x_i}{b_i+1}F(s),\qquad
 C_{\rm new}=\frac{x_i}{b_i+1}C.
\]
The definitions are linear in $fF$, so at any common admissible shift
\[
 \ell_{\rm new}(f)=\frac1{b_i+1}\ell_{\rm old}((s+x_i)f),
 \qquad Y_{\rm new}=Y_w(I+T_w/x_i).
\]
For a lower decrease $y_i\mapsto y_i-1$, with $y_i\ge2$, the same
calculation replaces $x_i$ by $y_i-1$.  A lower increase is its inverse
at the target, whenever that inverse exists.  The admissible shift
interval expands under an upper increase or lower decrease, so the two
intervals have a common element; shift independence then removes the
choice.

For the negative upper parameter, let the subscript on $\ell_m$ record
$m$ while $F,C,\eta$ stay fixed.  When $m\ge1$,
\[
 \omega_j(j-m)=-m(-1)^j\binom{m-1}{j},\qquad
 \sum_{j=0}^m\omega_j=0.
\]
In a divided difference the identity
\[
 \frac{(s-m)f(s)F(s)-(j-m)f(j)F(j)}{s-j}
 =(j-m)\frac{f(s)F(s)-f(j)F(j)}{s-j}+f(s)F(s)
\]
therefore implies
\[
 S_{m,F}[(s-m)f]=-mS_{m-1,F}[f],\qquad
 \ell_m((s-m)f)=-m\ell_{m-1}(f).
\]
Consequently
\[
 Y_{m-1}=Y_m(I-T_m/m),\qquad
 Y_m=Y_{m-1}(I-T_m/m)^{-1}
\]
when the inverse is defined.  This is exactly the native decrease
$x_3=-(m-1)\mapsto-m$, evaluated at the target.

\paragraph{6. Identification with the actual row state.}
At $(1,1,1,0;1,1,1)$, one has $F=C=1$, $m=\eta=0$ and
\[
 S_{f=1}=0,\quad S_{f=s}=1,\quad S_{f=s^2}=s,\qquad
 J_0(1)=1,\quad J_0(s)=-\tfrac12.
\]
Thus $Y_w=\bigl(\begin{smallmatrix}0&2&-1\\2&0&0\end{smallmatrix}\bigr)=R$.
Induction along any regular admissible path from this initial vertex,
using the generator equalities just established, gives $Y_w=RW$.
In particular, the formula applies along the stated four-block paths
provided all intermediate vertices satisfy the displayed domain and
inverse conditions.  The argument is uniform in the integer parameters;
the previously recorded finite checks are supporting diagnostics, not
its logical justification.  This algebraic identification does not
establish the additional prime-shell divisibility or any analytic
irrationality-measure estimate.


\paragraph{7. Domain check for the present trajectory.}
For $t=3$ or $4$, use the four blocks
\[
 (1,1,1,0;0,0,0)^{tN},\quad
 (0,0,0,0;2,1,3)^N,\quad
 (10-t,9-t,11-t,0;0,0,0)^N,\quad
 (0,0,0,-1;0,0,0)^{11N},
\]
where the superscript denotes the number of block steps and every elementary
move is evaluated at its current vertex.  During the lower-parameter block,
every old lower value is at most $3N$, whereas the three positive upper
values are $tN+1\ge3N+1$.  Hence the lower inverses are nonsingular.
All positive upper values are at least every lower value throughout the
four blocks, so $b_i\ge0$ and the shift interval is nonempty even at
intermediate coordinate moves.  Also
\[
 a_3=m-\sum_i b_i-3.
\]
It is negative in the first three blocks, where $m=0$, and in the last
block, where $\sum_i b_i=24N$ and $0\le m\le11N$.
For the final inverse at $m\ge1$, its determinant is
\[
 \det(I-T_m/m)=\frac{Q_m(m)}{a_3m^3}
 =\frac{\prod_i(m+y_i-1)}{a_3m^2}\ne0.
\]
The endpoint is therefore covered by the preceding uniform induction:
\[
 w=(10N+1,9N+1,11N+1,-11N;2N+1,N+1,3N+1),
 \quad 2N\le\eta\le9N.
\]
Here $F(s)=\binom{s+10N}{8N}\binom{s+9N}{8N}
\binom{s+11N}{8N}$.  This establishes the finite formula for the two
specified native rows on this trajectory, independently of any additional
prime-shell argument.


\section{Baseline arithmetic}\label{app:baseline}

% Independent audit: this lemma is conditional only on the finite CMF readout identity.
\paragraph{A baseline arithmetic lemma, with no extra-shell assumption.}
Let $N\ge1$, $m=11N$, $r\in\{0,1,2\}$, $E_r=13N+r$, and
\[
 F(s)=\binom{s+10N}{8N}\binom{s+9N}{8N}\binom{s+11N}{8N},\qquad
 c_j=(-1)^j\binom{m}{j}F(j).
\]
Put
\[
 S_r(s)=\sum_{j=0}^m(-1)^j\binom mj
              \frac{s^rF(s)-j^rF(j)}{s-j},\qquad
 J_\eta(P)=\sum_{t=0}^{\deg P}\frac{(-1)^t}{t+1}
                 \Delta^t P(-\eta).
\]
As usual $0^0=1$. The finite CMF readout identity established in Appendix~\ref{app:finite} is
\[
 C_NRW(N)e_r=2\binom{\sum_{j=0}^m j^r c_jH_{j+2N}^{(2)}+J_{2N}(S_r)}
                         {\sum_{j=0}^m j^r c_j}.
 \tag{F}
\]
The argument below proves arithmetic statements about the displayed finite
sums; its application to the matrix walk uses (F).

Set $D_{N,r}=d_{12N}d_{E_r}$ and define $\rho,\phi,G_N$ as defined in the main text. Then
\[
 \frac{D_{N,r}}{G_N}\,C_NRW(N)e_r\in\mathbb Z^2.
 \tag{B}
\]

\emph{Denominator calculation.}
If $f(s)=\binom{s+u}{b}$ and $j\in\mathbb Z$, then
$d_b(f(s)-f(j))/(s-j)$ is integer-valued. Indeed, the degree is at most
$b-1$, and its values at $j+1,\ldots,j+b$ are integers, since their
nonzero denominators divide $d_b$. Newton interpolation at these
consecutive integers proves the assertion at every integer.
A telescoping product difference now implies that $d_{8N}S_r$ is
integer-valued. The additional factor $s^r$ creates no denominator.
Expanding each divided difference into monomials and using
$\sum_j(-1)^j\binom mjj^t=0$ for $t<m$ gives
$\deg S_r\le24N-m+r-1=E_r-1$.
Thus $d_{8N}d_{E_r}J_{2N}(S_r)$ is an integer, and so is
$D_{N,r}J_{2N}(S_r)$.

For the harmonic term, fix a prime $p$ and a summand $j$.
If the highest power $P=p^a\le j+2N$ is at most $12N$, then
$v_p(D_{N,r})\ge2a$, which clears every reciprocal square in that
summand. Otherwise $12N<P\le j+2N\le13N$.
We have $p^{a-1}\le P/2\le13N/2<12N$, so
$v_p(D_{N,r})\ge2a-1$.
Also $j\ge P-2N$ and
\[
 8N<P,\qquad j+N\le12N<P,\qquad
 P\le j+9N\le20N<2P.
\]
The carry at the power $P$ in $\binom{j+9N}{8N}$ contributes a factor
$p$. All other prime-power contributions to its valuation are nonnegative.
Thus $p\mid c_j$, providing the missing factor. This proves that both
entries of $D_{N,r}C_NRW(N)e_r$ are integers.

\emph{Common divisor.}
Fix a prime $\sqrt{50N}<p\le12N$, and put $h=\phi(N/p)$.
All factorial arguments in $c_j$ are at most $22N<p^2$, hence
\[
 v_p(c_j)=\rho(N/p,j/p)\ge h.
 \tag{C}
\]
The minimum over $0\le y<1$ suffices because $\rho$ is periodic in $y$.
Use the sign-correct partial fraction identity
\[
 \sum_{j=0}^m\frac{(-1)^j\binom mj}{s-j}
   =\frac{(-1)^m m!}{\prod_{j=0}^m(s-j)}
\]
to write
\[
 S_r(s)=s^rF(s)\frac{(-1)^m m!}{\prod_{j=0}^m(s-j)}
                -\sum_{j=0}^m\frac{j^rc_j}{s-j}.
 \tag{P}
\]
Evaluate at the $E_r$ consecutive nodes $s=-1,-2,\ldots,-E_r$.
All nonzero integer factors that occur have absolute value at most
$E_r+m=24N+r<50N<p^2$.
If $F(s)\ne0$, counting signed consecutive factors gives exactly
\[
 v_p\!\left(F(s)\frac{m!}{\prod_{j=0}^m(s-j)}\right)
       =\rho(N/p,s/p)-1\ge h-1.
\]
For clarity, the factor $m!/\prod(s-j)$ contributes
$\lfloor m/p\rfloor-\lfloor(m-s)/p\rfloor+
\lfloor(-s-1)/p\rfloor$; for integer $s$ this is the first carry in
$\rho$ minus one. Each nonzero binomial polynomial factor is counted
by the corresponding difference of floors, also for negative arguments.
If $F(s)=0$, the first term of (P) vanishes.
Every other term of (P) has valuation at least $h-1$ by (C), because
$0<|s-j|<p^2$. Therefore $v_p(S_r(s))\ge h-1$ at all these nodes.

Since $d_{12N}S_r$ is integer-valued, its Newton coefficients at the
consecutive nodes are integers. Since $v_p(d_{12N})=1$, they are all
multiples of $p^h$. Translating an integer-valued binomial basis by an
integer preserves this divisibility, so the same holds for the
coefficients at $-2N$. The weights $1/(t+1)$ in $J_{2N}$ are cleared
by $d_{E_r}$ without losing $p^h$.
For the harmonic sum, $D_{N,r}$ contributes $p^2$,
all harmonic indices are less than $p^2$, and (C) supplies $p^h$.
The second row is a sum of the same integer coefficients.
Consequently $p^h$ divides both normalized entries.

It remains to check $12N<p\le13N$. Here
$x=N/p\in[1/13,1/12)$, and $\rho(x,0)=0$:
$\lfloor11x\rfloor=\lfloor10x\rfloor=\lfloor9x\rfloor=
\lfloor8x\rfloor=\lfloor3x\rfloor=\lfloor2x\rfloor=
\lfloor x\rfloor=0$. Every binomial carry is nonnegative, so
$\phi(x)=0$. Such primes contribute nothing to $G_N$.
Together with the established identity (F), this completes the proof of (B).



\section{Integral normalizations at the block boundaries}\label{app:blocks}

\paragraph{One complete multiplication route with matching normalizations.}
Use the four directions
\[
 (1,1,1,0;0,0,0),\quad(0,0,0,0;2,1,3),\quad
 (7,6,8,0;0,0,0),\quad(0,0,0,-1;0,0,0)
\]
for $3N,N,N,11N$ steps, respectively, with coordinate moves within
each step in the displayed order. The endpoint direction is
$(10,9,11,-11;2,1,3)$.
The completed-step vertices are
\[
\begin{aligned}
 w^{(1)}_k&=(k+1,k+1,k+1,0;1,1,1), &&0\le k\le3N,\\
 w^{(2)}_k&=(3N+1,3N+1,3N+1,0;2k+1,k+1,3k+1),&&0\le k\le N,\\
 w^{(3)}_k&=(3N+7k+1,3N+6k+1,3N+8k+1,0;
                2N+1,N+1,3N+1),&&0\le k\le N,\\
 w^{(4)}_k&=(10N+1,9N+1,11N+1,-k;2N+1,N+1,3N+1),&&0\le k\le11N.
\end{aligned}
\]
Every required lower inverse exists: its old lower parameter is at
most $3N$, strictly below all three positive uppers. Each required
upper decrease has a negative target parameter.
The coefficient $-a_3=\sum_{i=0}^3 x_i-\sum_{j=0}^2(y_j-1)$ is positive at every
intermediate vertex: it starts at three; increasing the positive
uppers increases it, the complete lower block leaves it at $3N+3$,
and the final negative-upper block leaves it at $13N+3$.
At completed steps its values are
$3k+3$, $9N-6k+3$, $3N+21k+3$, and $24N-k+3$.
Thus the described route has no singular generator.

Let $Y_i(k)$ be the actual two rows transported from $R$ along this
route, so $Y_i(k+1)=Y_i(k)B_i(k)$ and $Y_{i+1}(0)=Y_i(r_i)$.
At each vertex use
\[
 F(s)=\prod_{a=0}^2\binom{s+x_a-1}{x_a-y_a},\qquad C=F(0),
\]
and the finite readout formula (F), with the current $m=-x_3$ and
shift $\eta$. Write $B=\max(x_a-y_a)$, $E=\sum(x_a-y_a)-m+2$,
and $H=m+\eta$. The polynomial contribution is cleared by $d_Bd_E$.
If both lcm indices dominate $H$, they also clear the harmonic
contribution. The following choices therefore keep all six entries
integral and match at every boundary:
\[
\begin{array}{c|c|c|c|c|c}
\text{block}&m&\eta&B&E&L_i(k)\\\hline
1&0&0&k&3k+2&d_kd_{3k+2}\ (k>0)\\
2&0&2k&3N-k&9N-6k+2&d_{3N}d_{9N-6k+2}\\
3&0&2N&2N+6k&3N+21k+2&d_{\max(3N,2N+6k)}d_{3N+21k+2}\\
4&k&2N&8N&24N-k+2&d_{\max(8N,\min(12N,k+2N))}d_{24N-k+2}
\end{array}
\]
Set $L_1(0)=1$. In block 1 the harmonic term is zero.
In block 2 both indices dominate $H=2k$, and the first also
bounds $B=3N-k$.
In block 3 the three widths are
$N+7k,2N+6k,8k$; their maximum is $2N+6k$ for $0\le k\le N$.
Both lcm indices dominate $H=2N$.

In block 4 put $H=k+2N$ and $E=24N-k+2$.
We have $E\ge H$, and the first lcm index dominates the binomial
width $8N$. If $H\le12N$, both lcm indices dominate $H$.
Otherwise a missing reciprocal-square factor can occur only for
a prime power $P$ with $12N<P\le H\le13N$.
The preceding harmonic carry argument applies without alteration:
for a contributing $j\le k$ with $j+2N\ge P$,
$8N<P$, $j+N\le12N<P$, and
$P\le j+9N\le20N<2P$.
The binomial $\binom{j+9N}{8N}$ supplies the missing prime factor.
Thus $Z_i(k)=L_i(k)C_i(k)Y_i(k)$ is integral in every block.

The $C$ factors are evaluated from the same vertex, so they agree
at each boundary. The lcm factors there are, successively,
\[
 d_{3N}d_{9N+2},\qquad d_{3N}d_{3N+2},\qquad
 d_{8N}d_{24N+2}.
\]
Therefore $Z_{i+1}(0)=Z_i(r_i)$, with $Z_1(0)=R$.
The final normalization is $C_Nd_{12N}d_{13N+2}$ for all six
entries together. For one chosen column $r$, the same proof uses
its degree bound $E_r=\sum(x_a-y_a)-m+r$, replacing the final
second index by $13N+r$; at the initial point retain the explicit
normalization one. This gives the sharper columnwise endpoint
$C_Nd_{12N}d_{13N+r}$.

The shifts used above are admissible: each lies between the median
of the three lower offsets and the minimum positive-upper offset.
They are $0$, $2k$, $2N$, and $2N$ in the four blocks, respectively.
All binomial widths are nonnegative. These conditions are exactly
those used by the finite readout identity and its shift invariance.


\section{All six prime intervals}\label{app:all-shells}

\paragraph{All six extra prime shells: a uniform finite-field argument.}
This argument uses the established finite CMF readout, its admissible-shift
identity, and the baseline common-divisor lemma. It proves the full
additional divisor $H_{N,r}$, rather than extrapolating finite-depth tests.
Fix $r\in\{0,1,2\}$ and a prime $p$ with $p^2>50N$ lying in one of the
six column-corrected shells. Put
\[
 a=\lfloor8N/p\rfloor,\quad b=8N-ap,\qquad
 M=\lfloor11N/p\rfloor,\quad d=11N-Mp,\qquad h=\phi(N/p).
\]
The integer $h$ is the minimum carry count. The shell bounds and
$p^2>50N$ imply $p>50N/p>6$, so $p$ is odd; also
$a\le8N/p<p$ and $M\le11N/p<p$; also $a+1<p$. Work over $\mathbb F_p$ and
write
\[
 f_u(X)=\binom{X+uN}{b}\quad(u=9,10,11),\qquad
 f= f_9f_{10}f_{11},\quad g(X)=\prod_{j=0}^d(X-j),\quad z(X)=X^p-X.
\]
The binomials here are polynomials of degree $b<p$; their factorial
denominators are units.
For a residue $k\in\{0,\ldots,p-1\}$ define
\[
 e(k)=\mathbf1_{k>d},\quad
 \lambda_u(k)=\lfloor(k+uN)/p\rfloor,\quad
 \delta_u(k)=\mathbf1_{(k+uN)\bmod p<b},\qquad
 \chi(k)=e(k)+\sum_u\delta_u(k).
\]
Then $\chi(k)=\rho(N/p,k/p)\ge h$.
Each $f_u$ has a simple zero at $k$ exactly when $\delta_u(k)=1$.
Consequently
\[
 \mathcal Q(X)=\frac{f(X)}{z(X)^{h-1}g(X)}\in\mathbb F_p[X].
 \tag{Q}
\]
When $h=0$, this means $\mathcal Q=fz/g$, a polynomial because $g\mid z$.
For $h\ge1$, the multiplicity of each root $k$ in the denominator
is $h$ for $k\le d$ and $h-1$ for $k>d$, no larger than its
multiplicity in $f$.
Also $\mathcal Q(k)=0$ whenever $\chi(k)>h$.

\emph{The finite chamber table.}
Write $x=N/p$. On each open interval in the following table, the
minimum carry is the displayed $h$, and every residue with
$\chi(k)=h$ has the same displayed triple $(e,\lambda,\delta)$.
The triples are ordered by $u=9,10,11$.
\[
\begin{array}{c|cc|c|c|c|c|c}
 x\text{ interval}&a&M&h&e&\lambda&\delta&K=3a-M+h\\\hline
(1/8,2/13)&1&1&0&0&(1,1,1)&(0,0,0)&2\\
(2/9,3/13)&1&2&2&0&(2,2,2)&(1,1,0)&3\\
(3/10,4/13)&2&3&1&0&(2,3,3)&(0,1,0)&4\\
(3/8,5/13)&3&4&0&0&(3,3,4)&(0,0,0)&5\\
(4/9,5/11)&3&4&1&0&(4,5,5)&(0,1,0)&6\\
(5/11,6/13)&3&5&2&1&(4,5,5)&(0,1,0)&6\\
(8/9,10/11)&7&9&0&0&(8,9,10)&(0,0,0)&12\\
(10/11,12/13)&7&10&1&1&(8,9,10)&(0,0,0)&12
\end{array}
\tag{T}
\]
These are elementary floor inequalities, and can be checked by splitting
$0\le y<1$ at the moving boundaries
$\{cx\}$ for
$c\in\{11,-9,-10,-11,-1,-2,-3\}$.
Their ordering changes only when $cx$ or $(c-c')x$ is an integer.
The accompanying exact rational arrangement certificate checks every
open chamber and every interior boundary; it comprises 390 constant
floor cells, not a sampling of integer depths.
The omitted splitting values $x=5/11,10/11$ cannot occur for an eligible
prime: they force $(N,p)=(5,11)$ or $(10,11)$ after reduction, and
both violate $p^2>50N$.

For the selected column, its shell condition is $13N+r<Kp$.
Thus
\[
 \deg(X^r\mathcal Q)=13N+r-(K-1)p-1<p-1,
 \qquad\sum_{k\in\mathbb F_p}k^r\mathcal Q(k)=0.
 \tag{V}
\]
The last identity is the finite-field sum of a polynomial of degree
less than $p-1$, including its constant term.

\emph{Local factorization of the exact rational kernel.}
Let
\[
 F(s)=\prod_{u=9,10,11}\binom{s+uN}{8N},\qquad
 L(s)=\frac{(-1)^{11N}(11N)!}{\prod_{j=0}^{11N}(s-j)},\qquad
 R(s)=F(s)L(s).
\]
On the support $\chi(k)=h$, take $e,\lambda,\delta$ from (T) and set
\[
 P(V)=\prod_u (a+1)^{\delta_u}
                   \binom{\lambda_u+V}{a+\delta_u},\quad
 D_e(V)=\prod_{l=0}^{M-e}(V-l),\quad
 U(V)=\frac{P(V)}{D_e(V)},\quad
 \kappa=(-1)^{d-M+h-1}M!d!.
\]
These are independent of $k$. The factors in $D_e$ are distinct because
$M<p$. In the local coordinate $s=k+pV$, direct separation of the
factors divisible by $p$ gives
\[
 p^{1-h}R(k+pV)\equiv\kappa \mathcal Q(k)U(V)\pmod p.
 \tag{R}
\]
Here and below congruences of rational functions are taken away from
$V=0,\ldots,M-e$ before the removable factors are cancelled.
Here are the details of the factor count, including the case $e=1$.
For each binomial, if $\delta_u=0$ its normalized leading term is
$f_u(k)\binom{\lambda_u+V}{a}$.
If $\delta_u=1$, division by one $p$ gives
\[
 f'_u(k)(a+1)\binom{\lambda_u+V}{a+1}.
\]
This follows by splitting the $a$ or $a+1$ numerator factors divisible
by $p$; equivalently it is the first nonzero Lucas congruence.
It remains valid when an individual high-digit binomial vanishes.
For $L$, the $p$-factors in its denominator number $M+1-e$.
After their removal the remaining denominator is
$(-1)^M g'(k)$ if $e=0$, and $(-1)^M g(k)$ if $e=1$.
The numerator unit is $(-1)^dM!d!$.
Finally the product of the low-digit $f_u$ values or derivatives,
divided by $g'(k)$ or $g(k)$, is $(-1)^{h-1}\mathcal Q(k)$ by (Q).
These identities prove (R). For $\chi(k)>h$, the left side has zero
leading term and $\mathcal Q(k)=0$, so (R) continues to hold with zero right
side. All factor counts are exact at the first $p$-power level because
$p^2>50N$ bounds the relevant original factorial arguments.

Let $A_l=\operatorname{res}_{V=l}U(V)$, and let $B(V)$ be the polynomial
part of $U$:
\[
 U(V)=B(V)+\sum_{l=0}^{M-e}\frac{A_l}{V-l}.
\]
Set $c_j=0$ for indices outside $0\le j\le11N$. Taking residues in (R), using
$R(s)=S_0(s)+\sum_j c_j/(s-j)$, yields
\[
 p^{-h}c_{k+lp}\equiv\kappa \mathcal Q(k)A_l\pmod p.
 \tag{C'}
\]
For residue classes outside the minimal-carry support, every
$c_{k+lp}/p^h$ vanishes modulo $p$, including any extra endpoint term
$l=M$ when the supporting pattern has $e=1$.
For a general column, the exact kernel identity is
\[
 S_r(s)=s^rR(s)-\sum_{j=0}^{11N}\frac{j^rc_j}{s-j}.
\]
Multiplying by $p^{1-h}$ and subtracting the congruent-residue terms
from (R) now gives
\[
 p^{1-h}S_r(k+pV)
      \equiv\kappa k^r\mathcal Q(k)B(V)\pmod p.
 \tag{S'}
\]
Noncongruent terms in the residue sum vanish modulo $p$, since every
$c_j$ is divisible by $p^h$. This kernel argument is necessary when
$e=1$: termwise divided differences of $F$ would overlook contributing
noncongruent terms. Both sides of (S') are polynomials after subtraction,
so the identity includes the derivative nodes $V=l$.
More explicitly, the baseline interpolation argument established that
$T_r=p^{1-h}S_r$ is integer-valued over $\mathbb Z_{(p)}$.
Its degree is below $p^2$; substituting $s=k+pV$ into its Newton basis
therefore gives coefficients in $\mathbb Z_{(p)}$ (the numerator contains
at least $\lfloor t/p\rfloor$ factors of $p$ for the basis element of
degree $t$). Thus reduction and cancellation at these nodes are legitimate.

\emph{An admissible prime-dependent readout.}
Choose $\eta=Ap$, where $A=1$ if $p\ge2N$, and $A=2$ otherwise.
Every shell prime satisfies $2N\le\eta\le9N$, so this is an allowed
readout of the same matrix endpoint. For $j=k+lp$, the multiples of
$p$ not exceeding $j+\eta$ are exactly $p,\ldots,(l+A)p$.
The residue of $p^{2-h}$ times the harmonic part is therefore
\[
 \kappa\left(\sum_l A_l\sum_{t=1}^{l+A}t^{-2}\right)
           \sum_{k\in\mathbb F_p}k^r\mathcal Q(k)=0.
\]
The integers $1,\ldots,l+A$ are units modulo $p$, because the
harmonic indices are below $p^2$. Thus the harmonic part has valuation
at least $h-1$.

For the polynomial part, write $T_r=p^{1-h}S_r$.
The only weights in $J_\eta(T_r)$ possibly divisible by $p$ have index
$tp-1$, since $\deg S_r+1\le13N+r<p^2$.
At the Newton nodes $s=-Ap+up+k$, (S') gives
$T_r(s)\equiv\kappa k^r\mathcal Q(k)B(u-A)$.
Lucas's congruence gives
\[
 (-1)^{tp-1-(up+k)}\binom{tp-1}{up+k}
       \equiv(-1)^{t-1-u}\binom{t-1}{u}\pmod p.
\]
Consequently each $\Delta^{tp-1}T_r(-Ap)$ is a constant multiple of
$\sum_k k^r\mathcal Q(k)$, and is zero modulo $p$ by (V).
All weights in $J_\eta(T_r)$ are therefore cleared without losing a
factor of $p$: $J_\eta(T_r)\in\mathbb Z_{(p)}$.
The polynomial part, too, has valuation at least $h-1$.

The first entry of $C_NRW(N)e_r$ has valuation at least $h-1$;
the second is a sum of $c_j$ and has valuation at least $h$.
Since $v_p(d_{12N}d_{13N+r})=2$ and $v_p(G_N)=h$, both entries of
$C_Nd_{12N}d_{13N+r}RW(N)e_r/G_N$ are divisible by $p$.
Taking the product over all eligible primes in the six shells proves
\[
 \boxed{\frac{C_Nd_{12N}d_{13N+r}}{G_NH_{N,r}}
                    RW(N)e_r\in\mathbb Z^2.}
\]
The analytic section supplies the growth, error, and independence arguments needed to turn this arithmetic statement into an irrationality-measure bound.



\section{A reproducible rational enclosure for the arithmetic rate}\label{app:rate}
The moving breakpoints in the minimization over $y$ are
$\{sx\}$, for
$s\in\{11,-10,-9,-11,-2,-1,-3\}$, together with $0$ and $1$.
Their cyclic ordering can change only when $sx$ or $(s-t)x$ is an
integer. The separate floor $\lfloor8x\rfloor$ changes only when
$8x$ is integral. Thus it suffices to partition $[0,1]$ at every
$k/c$ for
\[
 c\in\{1,2,3,6,7,8,9,10,11,12,13,14,20,21,22\}.
\]
On each resulting open cell the minimum is constant: no breakpoint
crossing or individual floor change can occur. At a sample $x$ the
minimum is obtained by checking the $y$ breakpoints themselves and
one point in each gap. Values at the isolated $x$ boundaries do not
affect the integral. There are $94$ cells; the exact program computes
their heights $h_i$ and endpoints $a_i,b_i$. It gives
\[
 M:=\int_0^1\phi(x)\,dx=\frac{3517}{4095}.
\]
For a positive-height cell and integer $k\ge0$, its contribution is
$h_i((k+a_i)^{-1}-(k+b_i)^{-1})$.
The program sums the first $K=2000$ periods, rounding each positive
term down to a multiple of $10^{-15}$. There are $106000$ such terms;
their rounded sum is $12127998763267375/10^{15}$.
Periodicity and $k\le k+x\le k+1$ give the tail enclosure
\[
 \frac{M}{K+1}\le
 \int_K^\infty\frac{\phi(x)}{x^2}\,dx
 \le\frac{M}{K-1}.
\]
Adding the rounding error and the tail yields
\[
 \frac{159010385641752609841}{13110552000000000000}
 \le\gamma\le
 \frac{158851460347347083647}{13097448000000000000}.
\]
Numerically these endpoints are $12.12842797479104$ and
$12.128428404323277$. In particular
\[
 25-\gamma-\frac{233}{120}
 \le\frac{143297092558247390159}{13110552000000000000}
 <\frac{10929906}{1000000}.
\]
The final inequality is checked by integer cross multiplication in the supplied Python program. The finite partition and
directed summation are implemented with Python's exact
\texttt{Fraction} type and integer division; their uniform reduction
is the argument above, and their exhaustive evaluation is in
\texttt{rate\_certificate.py}. The proof uses these exact finite computations and the displayed uniform reductions.



\section{Short code accompanying the claims}
The full scripts, dependency pin, results, and commands are in
\texttt{code/}. The core RamanujanTools loop is reproduced here.
It checks finite native multiplication identities supporting the uniform generator argument.
\begin{lstlisting}[language=Python]
from ramanujantools import Position
from ramanujantools.cmf.pfq import pFq
import sympy as sp
from flint import fmpq, fmpq_mat

axes = sp.symbols('x0 x1 x2 x3 y0 y1 y2')
cmf = pFq(4, 3, 1)

def ordered_walk(N, t=3):
    w = [1, 1, 1, 0, 1, 1, 1]
    W = fmpq_mat([[1,0,0], [0,1,0], [0,0,1]])
    blocks = [(t*N, (1,1,1,0,0,0,0)),
              (N, (0,0,0,0,2,1,3)),
              (N, (10-t,9-t,11-t,0,0,0,0)),
              (11*N, (0,0,0,-1,0,0,0))]
    for count, direction in blocks:
        for _ in range(count):
            for i, amount in enumerate(direction):
                for _ in range(abs(amount)):
                    E = cmf.M(axes[i], amount >= 0).subs(
                        Position(dict(zip(axes, w))))
                    E = fmpq_mat([[fmpq(str(E[a,b]))
                          for b in range(3)] for a in range(3)])
                    W = W * E
                    w[i] += 1 if amount > 0 else -1
    return W

for N in (1, 2, 3):
    assert ordered_walk(N, 3) == ordered_walk(N, 4)
\end{lstlisting}

Run \texttt{python run\_checks.py} in the accompanying package to
execute the exact certificates and the finite diagnostics. The README
describes each program and its role. No proof-assistant files are included.


\begin{thebibliography}{9}
\bibitem{RhinViola1996}
G.~Rhin and C.~Viola,
\emph{On a permutation group related to $\zeta(2)$},
Acta Arithmetica \textbf{77} (1996), 23--56.
\href{https://doi.org/10.4064/aa-77-1-23-56}{doi:10.4064/aa-77-1-23-56}.
\href{https://www.math.ru.nl/~zudilin/zw/Rhin-Viola_AA_1996.pdf}{Full text}.

\bibitem{Zudilin2014}
W.~Zudilin,
\emph{Two hypergeometric tales and a new irrationality measure of $\zeta(2)$},
Annales math\'ematiques du Qu\'ebec \textbf{38} (2014), 101--117.
\href{https://arxiv.org/abs/1310.1526v3}{arXiv:1310.1526v3}.

\bibitem{David2023}
O.~David,
\emph{The conservative matrix field} (2023).
\href{https://arxiv.org/abs/2303.09318v3}{arXiv:2303.09318v3}.

\bibitem{Weinbaum2026}
S.~Weinbaum, E.~Leibtag, R.~Kalisch, M.~Shalyt and I.~Kaminer,
\emph{On Conservative Matrix Fields: Continuous Asymptotics and Arithmetic},
revised July 15, 2026.
\href{https://arxiv.org/abs/2507.08138v2}{arXiv:2507.08138v2}.
\end{thebibliography}


\end{document}
